\section{Image Classification}\label{sec:image-classification}

\subsection{Computer Vision and Digital Images}\label{sec:computer-vision-and-digital-images}

\subsubsection{What is Computer Vision?}\label{sec:what-is-computer-vision}

Humans continuously perform visual recognition: we recognize faces, identify objects, understand scenes, estimate distances, recognize actions, read text, notice anomalies. All of this happens almost effortlessly and automatically, without conscious thought. The question behind \definition{Computer Vision (CV)} is: \emph{\textbf{can we teach a computer to do the same?}} That is the entire purpose of the field.

\highspace
\begin{definitionbox}[: Computer Vision (CV)]
    \definition{Computer Vision (CV)} is an interdisciplinary scientific field that deals with how computers can be made to gain high-level understanding from digital images or videos.
\end{definitionbox}

\highspace
Let's make a few important observations about this definition:
\begin{itemize}
    \item ``\emph{Interdisciplinary}''. Computer Vision is not only Artificial Intelligence. It \textbf{combines ideas from several disciplines}, including Mathematics, Linear Algebra, Statistics, Signal Processing, Image Processing, Machine Learning, Deep Learning, Geometry, Optics, and Robotics.
    
    
    \item ``\emph{Digital images or videos}''. The input is always some visual data. For example, a JPEG image, a PNG image, a medical scan, a satellite image, frames from a camera, an entire video. However, notice something important. The \textbf{computer does not see objects}. It \hl{only receives numbers}.
    
    For example, an RGB image is simply: $I \in \mathbb{R}^{R \times C \times 3}$, where $R$ is the number of rows (height), $C$ is the number of columns (width), and $3$ is the number of channels (Red, Green, Blue). Later in this section, we'll see exactly how images are represented.
\end{itemize}
Thus, a computer starts from \textbf{low-level data} (numbers):
\begin{equation*}
    \begin{matrix}
        145 & 210 & 34 & \cdots \\
        144 & 208 & 36 & \cdots \\
        \vdots & \vdots & \vdots & \ddots
    \end{matrix}
\end{equation*}
These are just pixel intensities. Humans never think like this, instead we immediately understand ``\emph{there is a dog}'', or ``\emph{the dog is running}'', or ``\emph{the dog is looking left}''. These are \textbf{high-level concepts}. So Computer Vision tries to build a function:
\begin{equation*}
    f: \text{Pixels} \longrightarrow \text{Meaning}
\end{equation*}
Or more generally:
\begin{equation*}
    f(I) = \text{semantic information}
\end{equation*}
Where $I$ is the input image, and the output is some semantic information about the image. For example, the output could be a label (e.g., ``dog''), a bounding box (e.g., ``the dog is in this rectangle''), a segmentation mask (e.g., ``these pixels belong to the dog''), or even a caption (e.g., ``a dog is running in the park''). Thus, the whole challenge of Computer Vision is that \textbf{meaning is not explicitly present in the pixels}, the computer must \textbf{learn} it!
\begin{center}
    \begin{tabular}{@{} l | l @{}}
        \toprule
        Low-level information & High-level understanding \\
        \midrule
        Pixel intensities & Dog \\
        RGB values & Car \\
        Brightness & Person \\
        Colors & Bicycle \\
        Edges & ``A person riding a bike'' \\
        \bottomrule
    \end{tabular}
\end{center}

\highspace
\begin{flushleft}
    \textcolor{Green3}{\faIcon{question-circle} \textbf{Why is this difficult?}}
\end{flushleft}
Suppose we show a computer two images of a cat. To us, they look very similar. But to a computer they may look completely different because different colors, different lighting, different pose, different background, different size, etc. Although \textbf{semantically identical}, their pixel values may be almost unrelated. This is why Computer Vision is difficult. The mapping:
\begin{equation*}
    \text{pixels} \to \text{meaning}
\end{equation*}
Is highly \textbf{non-linear and extremely complex}.

\highspace
\textcolor{Green3}{\faIcon{question-circle} \textbf{Why has Computer Vision grown so much recently?}} Computer Vision has grown tremendously in the last decade, and is now a very active field of research. The main reasons are:
\begin{itemize}
    \item \textbf{Huge datasets}. \hl{Millions of labeled images became available.} Without large datasets, deep learning would not have achieved today's performance.
    \item \textbf{GPU computing}. Training a CNN involves billions of floating-point operations. \hl{Modern GPUs} made this computationally feasible.
    \item \textbf{Deep Learning}. Perhaps the biggest revolution. Before roughly 2012, most Computer Vision algorithms were \textbf{hand-designed}. Researchers manually created features such as: edges, corners, textures, gradients, SIFT, HOG, etc. And then fed them into classifiers. Today, \hl{Deep Neural Networks learn these features automatically} from data. This dramatically improved performance across almost every vision task.
\end{itemize}

\highspace
In summary \textbf{Computer Vision} aims to make computers understand the content of images and videos. The input is \textbf{low-level numerical data} (pixels), while the output is \textbf{high-level semantic information} (objects, scenes, actions). The central challenge is learning the mapping:
\begin{equation*}
    \text{Pixels} \to \text{Meaning}
\end{equation*}
Modern Computer Vision has advanced rapidly thanks to \textbf{large datasets}, \textbf{GPU computing}, and especially \textbf{Deep Learning}.